\chapter{Particle, Transport, Tensor-Train, and Tensor-Network Filters}
\label{ch:bf-highdim-particle-transport-tensor}

\section{Chapter Claim}
\label{sec:bf-highdim-ptt-claim}

Particle, transport, and tensor methods are plausible high-dimensional research
directions only when their approximation errors are measured in the downstream
filtering or inference task.  They reduce effective complexity when structure is
present; they do not remove the curse of dimensionality by declaration.

\section{Assumptions}
\label{sec:bf-highdim-ptt-assumptions}

The particle-filter discussion assumes finite proposal and target densities on
the tested support.  Transport and flow statements assume a differentiable map
with a computable proposal density when exact-target correction is claimed.
Tensor-train and tensor-network statements assume that low-rank structure is an
empirical or model-derived hypothesis to be tested by rank, normalization,
positivity, covariance, value, score, and downstream diagnostics.

\section{Particle Degeneracy}
\label{sec:bf-highdim-particle-degeneracy}

The bootstrap particle filter represents the predictive or filtering law by a
weighted empirical measure.  In high-dimensional observation regimes, likelihood
weights can concentrate on a tiny fraction of particles.  Resampling prevents
weight collapse from persisting forever, but it does not create missing
posterior support.  Guided proposals, auxiliary particle filters, localization,
and flow proposals all try to move particles closer to the next posterior before
weighting.

\section{Flow And Transport Proposals}
\label{sec:bf-highdim-flow-proposals}

A deterministic transport $T_\phi$ can push reference samples toward a
filtering proposal.  If the exact target density is still required, the
importance weight must include the proposal density and Jacobian correction.
Thus a learned flow is a proposal or preconditioner until a target-preserving
argument is supplied.

\paragraph{Transport proposal with correction boundary.}
\label{par:bf-highdim-transport-proposal}
\begin{enumerate}[label=(\arabic*)]
  \item Draw base samples $z^{(i)}$ from a reference proposal.
  \item Set proposed particles $x^{(i)}=T_\phi(z^{(i)})$.
  \item Evaluate proposal density using the change-of-variables formula.
  \item Weight particles against the desired prediction/update target.
  \item Record ESS, map residuals, log-Jacobian stability, and finite weights.
\end{enumerate}

\section{Ensemble Transport}
\label{sec:bf-highdim-ensemble-transport}

Ensemble transport filters and smoothers construct couplings between forecast
and analysis ensembles.  They are attractive because localization and triangular
or decomposable maps can express conditional structure.  The approximation is
the finite ensemble plus the chosen map class.  A transport update should be
judged by posterior checks, not only by map-training loss.

\section{Tensor-Train Filtering}
\label{sec:bf-highdim-tt-filtering}

Tensor-train methods represent a high-dimensional density or operator by a
chain of low-rank cores.  Direct nonlinear filtering work uses tensor trains to
approximate PDE-style filtering recursions or sequential Bayesian updates.
Sequential tensor-train Bayesian learning is especially relevant because
conditional transports and tensor decompositions can be connected to
NeuTra-style preconditioning.

The mandatory diagnostics are:
\begin{itemize}
  \item TT rank growth over time;
  \item density normalization error;
  \item positivity or signed-density defects;
  \item filtering likelihood or moment error against a low-dimensional
    reference;
  \item gradient parity if the tensor object is used inside HMC;
  \item posterior distortion under the actual downstream computation.
\end{itemize}

\section{Tensor-Network Kalman Filters}
\label{sec:bf-highdim-tn-kalman}

Tensor-network Kalman filters can compress large lifted linear systems, for
example Volterra-style nonlinear system identification after feature lifting.
The square-root tensor-network Kalman lesson is important for BayesFilter:
rounding compressed covariance factors can damage positive definiteness unless
the representation preserves a square-root or otherwise enforces covariance
validity.

\section{Low-Rank Tensor Observation Compression}
\label{sec:bf-highdim-low-rank-observation}

The diffusion MRI low-rank tensor UKF example is relevant as a design pattern:
compress the nonlinear observation geometry before filtering.  For BayesFilter,
the analogous question is whether a DSGE or nonlinear SSM observation/transition
map has block, low-rank, tensor-train, or sparse-factor structure that survives
posterior diagnostics.

\section{Pilot Architecture}
\label{sec:bf-highdim-tt-pilot}

A BayesFilter tensor pilot should start outside production defaults:
\begin{enumerate}[label=(\arabic*)]
  \item choose a low-dimensional or block-structured nonlinear fixture;
  \item fit or construct a TT/low-rank surrogate for a density, likelihood, or
    observation map;
  \item test normalization, positivity, value parity, score parity, and rank
    stability;
  \item use the surrogate only as a proposal or diagnostic until downstream
    posterior checks pass.
\end{enumerate}

\section{Literature Matrix}
\label{sec:bf-highdim-ptt-literature}

\begin{tabular}{p{0.20\linewidth}p{0.23\linewidth}p{0.20\linewidth}p{0.27\linewidth}}
\toprule
Family & Representative source & Support class & BayesFilter use \\
\midrule
TT nonlinear filtering & Li, Wang, Yau, Zhang, TT nonlinear filtering,
arXiv:1908.04010 & primary metadata; technical reading required & Direct
candidate for density/PDE lane. \\
Sequential TT Bayesian learning & Zhao and Cui, JMLR 2024 & primary metadata;
technical reading required & Bridge to transport and parameter learning. \\
Tensor-network Kalman & Batselier, Chen, Wong, arXiv:1610.05434 & primary
metadata; technical reading required & Lifted-system compression caution. \\
Square-root tensor Kalman & Menzen, Kok, Batselier, arXiv:2409.03276 & primary
metadata; technical reading required & PSD/covariance safeguard. \\
Low-rank tensor UKF & NeuroImage 2023 tractography paper & primary metadata;
domain-specific & Observation-compression design pattern. \\
Ensemble transport & Spantini/Baptista/Marzouk transport filtering line &
source needed for technical claims & Candidate localized transport lane. \\
\bottomrule
\end{tabular}

\section{Unresolved Claim Register}
\label{sec:bf-highdim-ptt-unresolved}

\begin{itemize}
  \item Technical sections of several tensor/transport papers still need
    full-paper review before theorem-level claims.
  \item No tensor backend exists in BayesFilter from this phase.
  \item No particle/transport method is promoted to HMC readiness.
\end{itemize}

\section{What Is Not Concluded}
\label{sec:bf-highdim-ptt-nonclaims}

This chapter does not conclude that particle, flow, transport, tensor-train, or
tensor-network methods solve high-dimensional filtering.  It defines research
lanes and the diagnostics required before any method can be treated as more than
a candidate.
